\section{Mixed evidence and institutional comparison}
The survey describes reported behavior; the authored accounts describe lived transactions; experts identify implementation problems; primary documents establish selected formal arrangements. Some migrant accounts are outbound but the headline survey outcome concerns disposition after receipt. No participant linkage is assumed. Four overlapping qualitative positions and cross-cutting security concerns are not fitted classes. Their wording and reported study coverage await coauthor review.

The benchmark contrasts Bre-B domestic interoperation; Pix public scheme, settlement and directory; phased linkage of Yape and Plin; SINPE Móvil numbers linked to local accounts; and SPEI distinct from CoDi. Record date and unit for each source. Separate interoperability, identity/onboarding, usability, remittance integration, recourse, gender/vulnerability and coordination. No national score is reissued. A documented route is not a recovery rate, and domestic infrastructure is not evidence of migrant uptake or cross-border delivery.

\section{Reproducibility and release boundary}
The current plan is to consider disclosure-reviewed code, documentation and selected aggregate outputs. Individual-level inputs and respondent-level auxiliary outputs remain restricted. No standing external verification access is offered. Before considering public files, check manual-exclusion code and other sources for embedded identifiers, and assess small cells and disclosure by subtraction across the package.

Retained consent materials and the reported oral explanation are distinct. Do not claim external ethics approval, exemption, or public-data authorization. Exact declarations and institutional decisions remain pending. The custodian brief is unsent, no contained private-build route is verified, and E.5 and Appendix C2 remain on hold. Historical LCA stays in the archive unless its separate scientific decisions are resolved.
